\documentclass[11pt]{article}
\usepackage[margin=1.15in]{geometry}
\usepackage{amsmath,amssymb,booktabs,tabularx}
\usepackage{microtype,xcolor,mdframed}
\usepackage[colorlinks=true,linkcolor=blue!50!black,urlcolor=blue!50!black]{hyperref}
\newmdenv[backgroundcolor=blue!4,linecolor=blue!35,linewidth=0.8pt,
  skipabove=9pt,skipbelow=9pt,innerleftmargin=10pt,innerrightmargin=10pt]{keyidea}
\newcommand{\code}[1]{{\small\ttfamily #1}}
\setlength{\emergencystretch}{1.5em}
\title{A Collatz Diary That Cannot Happen:\\
\Large A Friendly Guide to the $1/\sqrt2$ Sturmian Exclusion}
\author{Companion to ``An Unconditional Sturmian Exclusion\\for the Collatz Map at Slope $1/\sqrt2$''\thanks{Written with Codex (OpenAI). The guide explains the proved result; the companion paper and Lean source give the precise statements and proofs.}}
\date{September 4, 2026}
\begin{document}
\maketitle

\section{The result in ordinary language}
Take a positive whole number. If it is even, halve it. If it is odd,
triple it, add one, and halve the result. This combines two steps of the
usual Collatz rule whenever the number is odd. For example,
\[
3\longrightarrow5\longrightarrow8\longrightarrow4
\longrightarrow2\longrightarrow1\longrightarrow2\longrightarrow1\longrightarrow\cdots.
\]
Keep a diary: write $1$ when the current number is odd and $0$ when it is
even. The diary above begins $11000101\ldots$.

Our new result rules out a different family of infinite diaries. Imagine
walking along a number line in steps of length
\[
\alpha=\frac1{\sqrt2}\approx0.70710678.
\]
Choose any starting position $\rho$. At each step, write $1$ if you cross
into the next unit interval and $0$ if you remain in the same one. Since
the step is between zero and one, you never cross more than one boundary.
Starting at zero, this diary begins
\[
\underbrace{0\ 1\ 1\ 0\ 1\ 1\ 0\ 1\ 1\ 1\ 0\ 1\ 1\ 0\ 1\ 1\ 1\ 0\ 1\ 1}_{\text{first twenty letters}}\ \cdots.
\]
These are called \emph{mechanical words}; at an irrational step length
they are examples of Sturmian words. The starting position $\rho$ is the
\emph{intercept}. It changes the diary, but not the step length.

\begin{keyidea}
\textbf{What has been proved.} No whole-number Collatz orbit has any of
these infinite diaries, whatever intercept you choose. An orbit cannot
acquire one after an arbitrary finite beginning either. Lean has checked
the proof, including the approximation and size estimates.
\end{keyidea}
This excludes a particular structured family. It does not prove that
every Collatz orbit reaches $1$.

\newpage
\section{Why this family was worth attacking}
The fraction of ones in these diaries approaches $1/\sqrt2$, or about
$70.7\%$. In a rough multiplication-only picture, an odd shortcut step
contributes a factor of $3/2$ and an even step a factor of $1/2$. Over
$t$ steps containing $j$ ones, the multiplier is $3^j/2^t$. It changes
from shrinking to growing around $j/t=\log2/\log3\approx63.1\%$.

Our family is above that threshold. The usual shrinking intuition does
not rule it out. The proof must use the order of the zeros and ones,
and it must also keep the positive correction caused by adding one.
We do not assume that every possible runaway orbit has a fixed density
gap above the threshold.

\section{Matching diary pages force arithmetic agreement}
There is an exact fact behind the argument: two whole numbers with the
same first $m$ diary entries differ by a multiple of $2^m$.
For instance,
\begin{center}
\begin{tabular}{c@{\qquad}cccc}
\toprule
Start & Step 0 & Step 1 & Step 2 & Step 3\\
\midrule
$5$ & $5$ & $8$ & $4$ & $2$\\
$13$ & $13$ & $20$ & $10$ & $5$\\
\midrule
First three parities & $1$ & $0$ & $0$ & ---\\
\bottomrule
\end{tabular}
\end{center}
The first three diary entries agree, and $13-5=8=2^3$.
For distinct starts, agreement on a very long diary prefix requires a
very large difference between the starts.

Now apply this fact to two places in \emph{one} orbit. If the diary
repeats a block of length $q$ for a long stretch, the number at the start
of the stretch and the number $q$ steps later have many matching diary
entries. Either those two numbers are equal, making the future orbit
periodic, or the matching stretch forces a large arithmetic difference.

\begin{keyidea}
\textbf{The mechanism.} Find a long repeated stretch, make sure it appears
early enough, and compare two bounds: how large the matching diary says
the orbit must be, and how large the Collatz rule allows it to be.
\end{keyidea}
The words here never settle into a repeating pattern: their limiting
fraction of ones is irrational, whereas an eventually repeating diary
would have a rational limiting fraction. They nevertheless contain
long finite repetitions. The location of those repetitions matters:
a repetition found much later in the diary may give the orbit too much
time to grow. That is why the proof needs quantitative approximation.

\newpage
\section{A simple recurrence locates the repetitions}
The fractions
\[
\frac23,\quad\frac{12}{17},\quad\frac{70}{99},\quad
\frac{408}{577},\quad\ldots
\]
approach $1/\sqrt2$ from below. Their numerators $p$ and denominators $q$
come from the integer rule
\[
(q,p)\longmapsto(3q+4p,\ 2q+3p),\qquad(q,p)=(3,2)\text{ initially}.
\]
Every pair satisfies $q^2-2p^2=1$. Thus $q$ steps of the number-line walk
almost equal the whole number $p$, with a small positive leftover
$\delta=q/\sqrt2-p$.

Looking only at fractional parts turns the number line into a circle.
Shifting the diary by $q$ steps moves the fractional position by
$\delta$. Most positions retain the same letter. Changes occur only
near the boundary, in a short interval of length $\delta$. We call
arrivals in that interval \emph{visits}.

The proof establishes two complementary facts: visits cannot occur too
close together, and one cannot wait too long for the next visit.
Together they locate a long stretch between visits, which produces the
repeated diary block.

\begin{center}
\begin{tabular}{rrrr}
\toprule
Period $q$ & Numerator $p$ & Separation $Q$ & Window $G$\\
\midrule
3 & 2 & 7 & 17\\
17 & 12 & 41 & 99\\
99 & 70 & 239 & 577\\
577 & 408 & 1393 & 3363\\
\bottomrule
\end{tabular}
\end{center}
Here visits are at least $Q$ steps apart, and every window of $G$ steps
contains a visit. These statements hold for \emph{every} intercept.
The table shows the first few levels; Lean proves the recurrence and
its consequences for all levels.

\subsection*{The useful simplification}
The earlier plan called for the full three-distance theorem about
rotations. We found that a slightly larger rational grid is enough.
An elementary integer identity says the grid reaches every grid position.
A bound on the approximation error then guarantees a visit to the short
interval. The larger window costs something in the size estimate, but
there is still enough room for the contradiction. This removes a major
formalization dependency from this particular case.

\newpage
\section{Why the two bounds eventually clash}
Fix any proposed starting number $N$, however large. The repetition
argument forces a minimum size for certain orbit values. The exact
Collatz formula gives an upper bound for the same values. Combining
these estimates, the proof needs to choose a sufficiently large
recurrence denominator $q$ so that
\[
24(3N+21q+1)\left(\frac{a^6}{128}\right)^q<1,
\qquad a=3^{1/\sqrt2}.
\]
The constants are conservative bounds, not numerical guesses. The
important fact is $a<20/9$, and consequently $a^6<128$. So the fraction
inside the parentheses is strictly less than one.

The factor $3N+21q+1$ grows only linearly with $q$. The powered fraction
shrinks geometrically. For each fixed $N$, the shrinking factor eventually
wins. The recurrence supplies arbitrarily large $q$, so a suitable level
always exists. At that level the two orbit-size bounds contradict each
other. No candidate $N$ can keep the proposed infinite diary.

\section{What Lean checked, and what it did not claim}
Lean checked the integer recurrence, its approximation properties, the
visit argument for every intercept, the orbit estimates, and the final
limit argument. The theorem does not assume that the three-distance
theorem has been formalized, that a finite search has covered all starts,
or that an unproved statement about Collatz is available as an axiom.
Its dependencies are Lean's three standard foundational axioms.

\begin{keyidea}
\textbf{The remaining gap.} We have crossed one whole family off the list
of possible diaries. We have not shown that every hypothetical
counterexample must keep a diary in that family. Other Sturmian slopes,
less structured itineraries, and the general nontrivial-cycle problem
remain open.
\end{keyidea}
Related work from the same pass gives exact formulas for swapping nearby
diary entries and for organizing finite survivor sets. The paper records
these as tools, not as a claim that every orbit is now covered. It also
corrects an experimental size estimate that had omitted the positive
correction from Collatz's ``add one'' step.

\subsection*{Where to read the proof}
The technical companion is \code{paper/silver.pdf}; its editable source is
\code{paper/silver.tex}. The headline Lean theorem is
\code{Collatz.Silver.no\_\allowbreak{}itinerary}, in
\code{lean/Collatz/SturmianSilver.lean}. The ongoing research ledger is
\code{analysis/RESEARCH\_\allowbreak{}STATUS.md}. The new result is proved in the
repository; no claim that it is the first such result in the literature
is made here.
\end{document}
